\documentclass[11pt]{article}
\usepackage{amsmath,amssymb,amsthm}
\usepackage{booktabs}
\usepackage{enumitem}
\usepackage{hyperref}
\usepackage[margin=1in]{geometry}

\hypersetup{colorlinks=true,linkcolor=black,citecolor=blue,urlcolor=blue}

\newtheorem{theorem}{Theorem}
\newtheorem{lemma}[theorem]{Lemma}
\newtheorem{corollary}[theorem]{Corollary}
\newtheorem{definition}{Definition}
\newtheorem{remark}{Remark}
\newtheorem{proposition}[theorem]{Proposition}

\title{Zoo Per-LLM Chain Soundness\\
{\large Deterministic Gradient Ordering, Attribution-Graph
Correctness, and Byte-Identical Reproducibility}}
\author{Zoo Labs Foundation\\
\texttt{research@zoo.ngo}}
\date{Spec: 2026-04-15 \\ Effective: 2026-04-20}

\begin{document}
\maketitle

\begin{abstract}
The Zoo per-LLM training chain records every gradient update applied
to a model into a Quasar-finalised chain. We prove three results:
(i)~deterministic gradient ordering via Block-STM-as-consensus;
(ii)~correctness of the on-chain attribution graph that maps
training contributions to contributors; (iii)~a reproducibility
theorem stating that any party with the same training data, the same
random seed, and the same chain head produces a byte-identical model
artifact. Inheritance from \cite{quasar-cert} provides PQ finality.
\end{abstract}

\tableofcontents

%% ==================================================================
\section{Background}
\label{sec:bg}
%% ==================================================================

Zoo's per-LLM training chain is a distinct chain per model. Each
block in the chain records:
\begin{itemize}[nosep]
  \item a training step index $t$;
  \item a batch hash $\mathsf{b}_t$ (Merkle root over input batch);
  \item a gradient delta hash $\mathsf{g}_t$;
  \item a contributor identity $\mathsf{c}_t$ (DID);
  \item a model state root $\mathsf{m}_t$ post-update.
\end{itemize}

The chain is implemented on Lux's Block-STM-as-consensus
infrastructure (\cite{dagevm-lean}) with the additional constraint
that gradient updates are commutative within a single block --- as in
synchronous SGD --- but causally ordered across blocks.

%% ==================================================================
\section{Deterministic gradient ordering}
\label{sec:order}
%% ==================================================================

\begin{definition}[Gradient batch]
\label{def:batch}
A \emph{gradient batch} for step $t$ is the multiset
$G_t = \{ g_t^{(c)} \mid c \in \mathcal{C}_t \}$ where
$\mathcal{C}_t$ is the set of contributors who submitted a gradient
for step $t$.
\end{definition}

\begin{definition}[Aggregation rule]
\label{def:agg}
The aggregation function is the average:
$\bar{g}_t = \frac{1}{|\mathcal{C}_t|} \sum_{c \in \mathcal{C}_t} g_t^{(c)}$.
\end{definition}

\begin{lemma}[Order independence within a batch]
\label{lem:order-indep}
Floating-point summation is non-associative, but with a fixed
deterministic order (lexicographic on contributor DID), the result
is unique. The Zoo chain enforces this order at the block-encoding
level.
\end{lemma}

\begin{proof}
The block encoding sorts contributors by DID byte order before
aggregation. With fixed order, IEEE 754 produces a deterministic
result.
\end{proof}

\begin{theorem}[Deterministic ordering]
\label{thm:det-order}
Let $b$ be a per-LLM chain block finalised by a QuasarCert. The
post-block model state root $\mathsf{m}_b$ is a deterministic
function of the pre-state $\mathsf{m}_{b-1}$ and the gradient batch
$G_b$.
\end{theorem}

\begin{proof}
By Lemma~\ref{lem:order-indep}, aggregation is deterministic.
Block-STM serialisability~\cite{dagevm-lean} ensures one global
ordering of blocks. The model update
$\theta_{b} = \theta_{b-1} - \eta \cdot \bar{g}_b$ is a deterministic
function of $\theta_{b-1}$ and $\bar{g}_b$.
\end{proof}

%% ==================================================================
\section{Attribution graph}
\label{sec:attr}
%% ==================================================================

\begin{definition}[Attribution graph]
\label{def:attr}
The \emph{attribution graph} of a per-LLM chain is the directed
acyclic graph $\mathcal{G}_M = (V, E)$ where:
\begin{itemize}[nosep]
  \item $V$ = contributors recorded in any block of the chain;
  \item $(c, c') \in E$ iff $c$ contributed at step $t$ and $c'$
        contributed at step $t' > t$ on the same model lineage.
\end{itemize}
Edges are weighted by the marginal loss reduction attributable to
$c$'s contribution at step $t$.
\end{definition}

\begin{theorem}[Attribution-graph integrity]
\label{thm:attr-integ}
If the per-LLM chain is finalised under QuasarCert, then the
attribution graph $\mathcal{G}_M$ derivable from the chain is
unique and tamper-evident: any change to any
$(\mathsf{b}_t, \mathsf{g}_t, \mathsf{c}_t)$ tuple invalidates the
state root $\mathsf{m}_t$ and hence the QuasarCert at block $t$.
\end{theorem}

\begin{proof}
The block encoding stores
$(\mathsf{b}_t, \mathsf{g}_t, \mathsf{c}_t)$ as part of the block's
Merkle leaf. Modifying any field changes the block hash and
therefore $\mathsf{subj}_b$. By Theorem~7.5 of
\cite{quasar-cert}, forging a new QuasarCert on a modified subject
is bounded by $\mathsf{negl}(\lambda)$.
\end{proof}

\begin{corollary}[Contributor compensation correctness]
\label{cor:comp}
A contributor's compensation share, computed as the path-sum
weight in $\mathcal{G}_M$, is a deterministic function of the
chain. If compensation is paid via the Liquidity Protocol
(2026-04-20 adoption, see \cite{zoo-adopts}), the payment leg is
atomic with the attribution record.
\end{corollary}

%% ==================================================================
\section{Byte-identical reproducibility}
\label{sec:repro}
%% ==================================================================

\begin{definition}[Reproducibility predicate]
\label{def:repro}
Given a per-LLM chain head $b^\star$ and a public initial state
$\mathsf{m}_0$, the chain is \emph{byte-identically reproducible}
iff: any party that:
\begin{enumerate}[nosep]
  \item starts from $\mathsf{m}_0$;
  \item replays each block $b = 1, \ldots, b^\star$ in chain order;
  \item applies the canonical aggregation
        (Lemma~\ref{lem:order-indep});
\end{enumerate}
produces the same byte-string final model
$\theta_{b^\star}$ as the original training run.
\end{definition}

\begin{theorem}[Reproducibility]
\label{thm:repro}
The Zoo per-LLM chain is byte-identically reproducible.
\end{theorem}

\begin{proof}
By Theorem~\ref{thm:det-order} the per-block update is deterministic
given the pre-state and the gradient batch. By
Theorem~\ref{thm:attr-integ} the gradient batches are
tamper-evidently recorded. By induction on block number, the final
state $\theta_{b^\star}$ depends only on $\mathsf{m}_0$ and the
recorded batches. Floating-point determinism is enforced by
canonical contributor ordering (Lemma~\ref{lem:order-indep}); the
training framework binds to a single FP implementation
(IEEE 754 binary32 / binary16) and a deterministic CUDA / ROCm
kernel set declared at chain genesis. Hence anyone can rebuild the
final model byte-for-byte.
\end{proof}

\begin{remark}[FP determinism caveat]
GPU non-determinism is real. Zoo per-LLM chains use the
NVIDIA / AMD ``deterministic ops'' subset and pin all kernel
versions in the chain genesis. Under these constraints the
proof goes through. If a GPU vendor introduces a non-deterministic
kernel, the chain genesis must be updated (a fork).
\end{remark}

%% ==================================================================
\section{Detailed proof: deterministic ordering}
\label{sec:detail-order}
%% ==================================================================

We expand Theorem~\ref{thm:det-order}.

\paragraph{Setup.}
At step $t$, the model state is $\theta_{t-1}$ (a vector of
floating-point weights) and the gradient batch is $G_t$ (a multiset
of contributor gradient vectors). We assume:
\begin{itemize}[nosep]
  \item gradient vectors have a fixed dimension $d$ (the model's
        parameter count);
  \item every $g^{(c)}$ is encoded as IEEE 754 binary32 (or
        binary16) in row-major order;
  \item the chain genesis pins the GPU/CPU kernel manifest
        $\mathrm{KernelManifest}$.
\end{itemize}

\paragraph{Ordering subroutine.}
The block builder sorts contributors lexicographically by DID byte
order; let $c_{(1)}, \ldots, c_{(k)}$ be the sorted list. The
aggregate is computed as a left-fold:
$\bar{g}_t = \frac{1}{k} \cdot
  ((((g^{(c_{(1)})} + g^{(c_{(2)})}) + g^{(c_{(3)})}) +
    \cdots) + g^{(c_{(k)})}).$
This left-fold order is unique given the sorted input.

\paragraph{Determinism of left-fold.}
IEEE 754 addition is non-associative, but for fixed left-fold order
the result is unique. The kernel manifest pins the FP rounding
mode (round-to-nearest-even) and the CUDA / ROCm fused-multiply-add
flag (off by default in deterministic mode).

\paragraph{Update step.}
$\theta_t = \theta_{t-1} - \eta \cdot \bar{g}_t$ where $\eta$ is the
learning-rate schedule (also pinned in the chain genesis as a
deterministic function of $t$).

\paragraph{Result.}
$\theta_t$ is a pure function of $(\theta_{t-1}, G_t,
\mathrm{KernelManifest}, \eta(t))$, all of which are recorded on
chain. Hence the post-block model state is deterministic, as
claimed in Theorem~\ref{thm:det-order}.

%% ==================================================================
\section{Per-block storage layout}
\label{sec:layout}
%% ==================================================================

\begin{table}[h]
\centering
\caption{Per-LLM chain block field layout.}
\begin{tabular}{lll}
\toprule
Field & Type & Size \\
\midrule
\texttt{chain\_id} (model registry id)   & \texttt{uint64}        & 8 B  \\
\texttt{step\_index} $t$                 & \texttt{uint64}        & 8 B  \\
\texttt{batch\_root} $\mathsf{b}_t$       & \texttt{bytes32}       & 32 B \\
\texttt{gradient\_root} $\mathsf{g}_t$    & \texttt{bytes32}       & 32 B \\
\texttt{contributors} (sorted DID list)  & \texttt{bytes[]}       & variable \\
\texttt{model\_root} $\mathsf{m}_t$       & \texttt{bytes32}       & 32 B \\
\texttt{quasar\_cert} $\mathcal{C}_t$     & cert tuple             & $\approx 8$ KB \\
\bottomrule
\end{tabular}
\end{table}

The contributor list is the canonical lexicographic-DID-byte order
(Lemma~\ref{lem:order-indep}); any non-canonical ordering
invalidates the deterministic-aggregation premise.

%% ==================================================================
\section{Reproducibility test plan}
\label{sec:test}
%% ==================================================================

For any per-LLM chain $\mathcal{M}$ a third party can:

\begin{enumerate}[nosep]
  \item snapshot the chain head $b^\star$;
  \item read the genesis $(\mathsf{m}_0,
        \mathrm{KernelManifest})$ from block 0;
  \item replay each block's gradient batch in order, using only
        the data accessible from the chain (batch hashes index
        into IPFS; gradient hashes index into the contributor's
        commitment);
  \item compute $\theta_{b^\star}$;
  \item verify $H(\theta_{b^\star}) = \mathsf{m}_{b^\star}$.
\end{enumerate}

Theorem~\ref{thm:repro} states that this test always succeeds,
modulo the FP-determinism caveat (kernel pinning).  In practice,
Zoo uses NVIDIA cuBLAS deterministic mode and AMD rocBLAS
deterministic kernel set as declared in
$\mathrm{KernelManifest}$.

%% ==================================================================
\section{Cross-references with Zoo ZIPs}
\label{sec:zip}
%% ==================================================================

\begin{itemize}[nosep]
  \item ZIP-0400 -- AI Assistant: defines the inference / training
        contract surface.
  \item ZIP-0401 -- Per-LLM chains: introduces this construction.
  \item ZIP-0405 -- Reproducible training: pins the deterministic
        kernel manifest format.
  \item ZIP-0418 -- Contributor compensation: defines the
        attribution-graph weight formula and the on-chain payout.
\end{itemize}

%% ==================================================================
\section{Composition with Liquidity Protocol}
\label{sec:lp-compose}
%% ==================================================================

Contributor compensation, when paid on-chain, is settled via the
Liquidity Protocol post 2026-04-20 (see \cite{zoo-adopts}).
Composition follows from the same argument as
\cite{lux-adopts} \S4: trade atomicity composes with Block-STM
serialisability; the per-LLM chain is itself a Block-STM-as-consensus
chain.

%% ==================================================================
\begin{thebibliography}{9}

\bibitem{quasar-cert}
Z.~Kelling.
\newblock QuasarCert Soundness, Quasar 3.0.
\newblock \texttt{luxfi/proofs/quasar-cert-soundness.tex}, 2025-12-15.

\bibitem{dagevm-lean}
Lux Industries.
\newblock \texttt{lean/Consensus/DAGEVM.lean}.

\bibitem{zoo-adopts}
Zoo Labs Foundation.
\newblock Zoo Adopts the Liquidity Protocol.
\newblock \texttt{zooai/proofs/zoo-adopts-liquidity-protocol.tex}, 2026-04-20.

\bibitem{lux-adopts}
Lux Industries.
\newblock Lux Adopts the Liquidity Protocol.
\newblock \texttt{luxfi/proofs/lux-adopts-liquidity-protocol.tex}, 2026-04-20.

\bibitem{zoo-2021}
A.~Worring, Z.~Kelling.
\newblock Zoo: NFT-Driven Conservation, October 2021.
\newblock \texttt{zooai/papers/zoo-2021-original-whitepaper/}.

\end{thebibliography}

\end{document}
